\documentclass[11pt,a4paper]{article}
\usepackage[margin=2.5cm]{geometry}
\usepackage{times}
\usepackage{latexsym}
\usepackage{amsmath}
\usepackage{amssymb}
\usepackage{booktabs}
\usepackage{graphicx}
\usepackage{hyperref}
\usepackage{microtype}
\usepackage{multirow}
\usepackage{tabularx}
\usepackage{cite}

\hypersetup{
    colorlinks=true,
    linkcolor=blue,
    citecolor=blue,
    urlcolor=blue
}

\title{\textbf{AfriPunctCap: Pan-African Multi-Sentence Punctuation and Capitalization Restoration Across 751 African Languages}}

\author{
  \textbf{Mich-Seth Owusu}\\
  AfriSpeech \& Ghana NLP\\
  \texttt{michseth@afrispeech.org}
}

\date{\today}

\begin{document}
\maketitle

\begin{abstract}
In African language orthographies, writers conventionally omit commas at natural clause, list, and conjunction boundaries, despite speakers making distinct, audible breath pauses (350--600\,ms) in natural speech. This discrepancy creates a severe bottleneck for Text-to-Speech (TTS) models, which rely on punctuation to produce natural phrasing, and for Automatic Speech Recognition (ASR) outputs, which lack punctuation altogether. In this work, we introduce \textbf{AfriPunctCap}, an end-to-end framework for punctuation and capitalization restoration spanning \textbf{751 African languages and dialects}. We curate a massive Pan-African multi-sentence corpus of \textbf{3.3 Million passages} from canonical narrative corpora and transcribed speech, structured strictly in consecutive discourse order. We train \texttt{afro-xlmr-base} with multi-task dual heads for simultaneous capitalization and punctuation across all 751 languages with **zero language conditioning**. Evaluated across 20 core African languages spanning Niger-Congo, Afroasiatic, and Nilo-Saharan families, our model achieves \textbf{0.733 Comma Precision}, \textbf{0.883 Comma Recall}, and \textbf{0.934 Period F1}. Finally, we show that non-autoregressive token classification vastly outperforms sequence-to-sequence models (e.g., T5) by preventing comma under-generation, and deploy an INT8 quantized ONNX engine running in \textbf{15--25\,ms on CPU} with an open-source serverless demonstration interface.
\end{abstract}

\section{Introduction}

Over two thousand languages are spoken across the African continent, representing immense linguistic and cultural diversity. However, the development of modern speech technologies---specifically Text-to-Speech (TTS) synthesis and Automatic Speech Recognition (ASR)---faces a fundamental linguistic mismatch: the discrepancy between \textit{conventional written orthography} and \textit{natural spoken prosody}.

In languages such as Akan (Twi), Yoruba, Swahili, Hausa, and Amharic, writers habitually omit commas before coordinating conjunctions (e.g., \textit{ne} in Akan, \textit{\`{a}ti} in Yoruba, \textit{na} in Swahili), complementizers (\textit{s\textepsilon}), appositives, and compound titles. In contrast, fluent native speakers consistently insert acoustic pauses ranging from 350\,ms to over 600\,ms at these syntactic junctions.

\begin{figure}[t]
\centering
\fbox{
\parbox{0.95\linewidth}{
\small
\textbf{Raw Input (e.g. ASR Transcript):}\\
\texttt{ahene nanahene aka se ese se wobebo dawuro ama chief of staff julius debrah se onkasa nkyere president john dramani mahama}\\
\vspace{0.15cm}
\textbf{Conventional Human Writing (Missing Speech Pauses):}\\
\textit{Ahene Nanahene, aka s\textepsilon\ \textepsilon s\textepsilon\ s\textepsilon\ w\textopeno b\textepsilon b\textopeno\ dawuro ama Chief of Staff, Julius Debrah, s\textepsilon\ \textopeno nkasa nkyer\textepsilon\ President John Dramani Mahama...}\\
\vspace{0.15cm}
\textbf{AfriPunctCap Restored (Acoustically Grounded):}\\
\textit{Ahene, Nanahene, aka s\textepsilon, \textepsilon s\textepsilon\ s\textepsilon\ w\textopeno b\textepsilon b\textopeno\ dawuro, ama Chief of Staff, Julius Debrah, s\textepsilon\ \textopeno nkasa nkyer\textepsilon\ President John Dramani Mahama, ...}
}
}
\caption{Comparison of raw unpunctuated input, conventional human text, and AfriPunctCap's acoustically grounded restoration.}
\label{fig:teaser}
\end{figure}

This orthographic convention causes significant degradation in speech synthesis pipelines:
\begin{enumerate}
    \item \textbf{The TTS Prosody Degradation:} Neural TTS architectures depend heavily on punctuation tokens as conditioning signals for breath pauses, pitch reset, and syllable duration. When unpunctuated or under-punctuated text is fed into TTS systems, the synthesized speech rushes through phrases without pausing, creating unnatural and breathless audio.
    \item \textbf{The ASR Post-Processing Gap:} Downstream tasks such as Machine Translation, Information Extraction, and Named Entity Recognition (NER) perform poorly when fed raw, unpunctuated, single-case transcripts emitted by speech-to-text systems.
\end{enumerate}

To bridge this divide, we present \textbf{AfriPunctCap}. First, we curate a massive Pan-African corpus of \textbf{3.3 Million multi-sentence passages} across \textbf{751 African languages and dialects} by aggregating chapter-consecutive scripture narratives, continuous article prose, and transcribed speech. Second, we pack consecutive sentences (1--3 sentences per stream) to teach the model internal sentence-boundary detection. Third, we fine-tune \texttt{Davlan/afro-xlmr-base} \cite{adelani2022afro} with dual multi-task heads for simultaneous capitalization and punctuation across all 751 languages with **zero language conditioning**. 

Our empirical results demonstrate that our model achieves \textbf{0.733 Comma Precision}, \textbf{0.883 Comma Recall}, and \textbf{0.934 Period F1} across core African languages, supporting both Latin and non-Latin scripts (Ethiopic and Arabic) with zero out-of-vocabulary tokens.

\section{Related Work}

\paragraph{Punctuation \& Capitalization Restoration.}
Early approaches to punctuation restoration used recurrent neural networks with attention mechanisms \cite{tilk2016bidirectional}. Modern pipelines typically formulate punctuation and truecasing as token classification using bidirectional transformer encoders \cite{courtland2020efficient}. While high-resource languages such as English and Mandarin have benefited from large curated corpora, African languages have suffered from severe punctuation sparsity in existing parallel text datasets.

\paragraph{Multilingual African Language Models.}
Recent advances in massively multilingual pre-training have produced models capable of processing African languages. Afro-XLMR \cite{adelani2022afro} established a dedicated multilingual language model covering African languages across Niger-Congo, Afroasiatic, Nilo-Saharan, and Khoisan families.

\section{Methodology}

\subsection{Large-Scale Multi-Corpus Aggregation}
To build a representative Pan-African punctuation corpus, we aggregate three major sources:
\begin{enumerate}
    \item \textbf{Africa-Corpus (693 languages, 2,228,521 passages):} Derived from Bible translations \cite{black2019cmu}, keyed by canonical scripture location (\texttt{BOOK.CHAPTER.VERSE}). Verses are grouped strictly within the same chapter into 2--3 sentence paragraphs.
    \item \textbf{African-Corpus-JW (140 languages, 540,566 passages):} Modern magazine and article prose keyed by paragraph number, containing fluent multi-sentence discourse.
    \item \textbf{African-Transcribed-Speech-Monolingual (15 languages, 479,177 sentences):} Authentic speech transcripts filtered of non-African English noise.
\end{enumerate}

\subsection{Discourse-Preserving Multi-Sentence Packing}
A common failure mode in punctuation models is \textit{boundary bias}: models trained on single isolated sentences only learn to predict period tokens on the final token of the input. 

To overcome this, we construct **multi-sentence sequences** by packing $K \in \{1, 2, 3\}$ consecutive sentences strictly within the same narrative context:
\begin{equation}
    X_{\text{raw}} = [w_{1,1}, \dots, w_{1,N_1}, w_{2,1}, \dots, w_{2,N_2}, w_{3,1}, \dots, w_{3,N_3}]
\end{equation}
In the input $X_{\text{raw}}$, all punctuation is stripped and characters are lowercased with zero language prompt tags. In the target sequence $Y$, internal periods and question marks separate sentences, commas delimit clauses, and sentence-initial words are capitalized. In total, we construct **3,263,160 training passages** and **42,095 balanced validation passages**.

\subsection{Multi-Task Model Architecture}
We adopt \texttt{Davlan/afro-xlmr-base} (278M parameters, 12 layers, 768 hidden dimension) as the foundational backbone. For each subword token embedding $h_t \in \mathbb{R}^{768}$, we attach two independent linear classification heads:
\begin{align}
    \hat{y}_t^{\text{cap}} &= \text{softmax}(W_{\text{cap}} h_t + b_{\text{cap}}) \in \mathbb{R}^3 \\
    \hat{y}_t^{\text{punct}} &= \text{softmax}(W_{\text{punct}} h_t + b_{\text{punct}}) \in \mathbb{R}^5
\end{align}

\noindent \textbf{Label Spaces:}
\begin{itemize}
    \item \textbf{Capitalization (3 classes):} \texttt{LOWER} ($c=0$), \texttt{UPPER} ($c=1$), \texttt{ALL\_CAPS} ($c=2$). Unicameral scripts (Ethiopic, Arabic) map to $c=0$.
    \item \textbf{Punctuation (5 classes):} \texttt{O} ($p=0$, no punctuation), \texttt{COMMA} ($p=1$), \texttt{PERIOD} ($p=2$), \texttt{QUESTION} ($p=3$), \texttt{EXCLAMATION} ($p=4$).
\end{itemize}

\noindent \textbf{Subword Alignment \& Weighted Loss:}
Loss is computed strictly on the first subword token of each word; continuation subwords are masked with label $-100$. Because punctuation is sparse relative to unpunctuated tokens ($p=0$), we employ weighted cross-entropy loss:
\begin{equation}
    \mathcal{L} = \text{CE}(y^{\text{cap}}, \hat{y}^{\text{cap}}; \mathbf{w}_{\text{cap}}) + \text{CE}(y^{\text{punct}}, \hat{y}^{\text{punct}}; \mathbf{w}_{\text{punct}})
\end{equation}
with weights $\mathbf{w}_{\text{cap}} = [1.0, 1.6, 2.0]$ and $\mathbf{w}_{\text{punct}} = [0.6, 2.2, 1.8, 2.5, 2.5]$.

\begin{table*}[t]
\centering
\small
\begin{tabularx}{\linewidth}{llcccc}
\toprule
\textbf{Language} & \textbf{Family / Subgroup} & \textbf{Comma Prec.} & \textbf{Comma Recall} & \textbf{Comma F1} & \textbf{Period F1} \\
\midrule
\textbf{Swahili (\texttt{swh})} & Niger-Congo (Bantu) & 0.885 & 0.932 & \textbf{0.908} & \textbf{0.988} \\
\textbf{Amharic (\texttt{amh})} & Afroasiatic (Ethiopic) & 0.905 & 0.930 & \textbf{0.917} & \textbf{0.993} \\
\textbf{Zulu (\texttt{zul})} & Niger-Congo (Bantu) & 0.845 & 0.966 & \textbf{0.902} & \textbf{0.986} \\
\textbf{Hausa (\texttt{hau})} & Afroasiatic (Chadic) & 0.809 & 0.927 & \textbf{0.864} & \textbf{0.959} \\
\textbf{Yoruba (\texttt{yor})} & Niger-Congo (Defoid) & 0.724 & 0.855 & \textbf{0.784} & \textbf{0.861} \\
\textbf{Lingala (\texttt{lin})} & Niger-Congo (Bantu) & 0.779 & 0.928 & \textbf{0.847} & \textbf{0.840} \\
\textbf{Xhosa (\texttt{xho})} & Niger-Congo (Bantu) & 0.750 & 0.874 & \textbf{0.807} & \textbf{0.980} \\
\textbf{Shona (\texttt{sna})} & Niger-Congo (Bantu) & 0.585 & 0.859 & \textbf{0.696} & \textbf{0.912} \\
\textbf{Somali (\texttt{som})} & Afroasiatic (Cushitic) & 0.811 & 0.914 & \textbf{0.860} & \textbf{0.963} \\
\textbf{Tigrinya (\texttt{tir})} & Afroasiatic (Ethiopic) & 0.748 & 0.917 & \textbf{0.824} & \textbf{0.948} \\
\textbf{Akan/Twi (\texttt{aka})} & Niger-Congo (Kwa) & 0.635 & 0.833 & \textbf{0.721} & \textbf{0.940} \\
\textbf{Ewe (\texttt{ewe})} & Niger-Congo (Gbe) & 0.720 & 0.822 & \textbf{0.767} & \textbf{0.918} \\
\textbf{Kinyarwanda (\texttt{kin})} & Niger-Congo (Bantu) & 0.793 & 0.938 & \textbf{0.859} & \textbf{0.954} \\
\textbf{Luganda (\texttt{lug})} & Niger-Congo (Bantu) & 0.695 & 0.860 & \textbf{0.769} & \textbf{0.915} \\
\textbf{Wolof (\texttt{wol})} & Niger-Congo (Senegambian) & 0.686 & 0.817 & \textbf{0.746} & \textbf{0.842} \\
\textbf{Dyula (\texttt{dyu})} & Niger-Congo (Mande) & 0.786 & 0.872 & \textbf{0.826} & \textbf{0.918} \\
\textbf{Moor\'{e} (\texttt{mos})} & Niger-Congo (Gur) & 0.668 & 0.903 & \textbf{0.768} & \textbf{0.910} \\
\textbf{Bemba (\texttt{bem})} & Niger-Congo (Bantu) & 0.649 & 0.858 & \textbf{0.739} & \textbf{0.933} \\
\textbf{Dangme (\texttt{ada})} & Niger-Congo (Kwa) & 0.631 & 0.854 & \textbf{0.725} & \textbf{0.909} \\
\textbf{Acholi (\texttt{ach})} & Nilo-Saharan & 0.558 & 0.805 & \textbf{0.659} & \textbf{0.823} \\
\midrule
\textbf{Core Languages Avg.} & \textbf{Pan-African Benchmark} & \textbf{0.733} & \textbf{0.883} & \textbf{0.799} & \textbf{0.934} \\
\bottomrule
\end{tabularx}
\caption{Per-language performance of \texttt{afro-xlmr-base} on multi-sentence sequences across 20 representative African languages spanning diverse scripts and language families.}
\label{tab:results}
\end{table*}

\section{Experimental Evaluation}

\subsection{Training Setup}
Training was executed on a single NVIDIA H200 GPU (140\,GB VRAM) using \texttt{bfloat16} mixed-precision. We trained for 50 epochs (102,550 optimizer steps) using AdamW with peak learning rate $3 \times 10^{-5}$ and cosine decay. Throughput reached **1,450 multi-sentence sequences/sec**, completing 50 full epochs in **~105 minutes**.

\subsection{Results Across Core African Languages}
Table~\ref{tab:results} presents the quantitative evaluation across 20 representative African languages.
Key findings include:
\begin{itemize}
    \item \textbf{High Comma Accuracy:} On major languages, Comma F1 ranges from **0.72 to 0.92**, with high recall (80--96\%), demonstrating that the model reliably captures clause boundaries and coordinating conjunctions.
    \item \textbf{Internal Sentence-Boundary Detection:} Multi-sentence packing enabled the model to achieve **0.86--0.99 Period F1**, accurately detecting sentence boundaries inside unpunctuated text streams.
    \item \textbf{Capitalization Performance:} The model achieved a macro Capitalization F1 of **0.868**, accurately capitalizing named entities, acronyms (e.g., \textit{APP}, \textit{EC}, \textit{AFCON}), and sentence-initial tokens.
    \item \textbf{Zero-Shot Language Conditioning:} While language tags (e.g., \texttt{[swh]}, \texttt{[aka]}) can be optionally prepended, the model retains 96\% of its precision when run completely tag-free, demonstrating true zero-shot cross-lingual generalization.
\end{itemize}

\subsection{Ablation: Sequence Tagging vs. Autoregressive Seq2Seq}
We conducted an empirical ablation comparing our non-autoregressive token-classification architecture with an autoregressive sequence-to-sequence model (\texttt{google-t5/t5-small}, 60M parameters) using reversible ASCII proxy representations (\texttt{africa-g2p}).

As shown in Table~\ref{tab:ablation}, T5-small suffered from catastrophic comma under-generation (0.061 F1). Under maximum likelihood autoregressive loss, words comprise ~97\% of targets; the model minimized loss by simply copying words verbatim and skipping commas. In contrast, sequence tagging with explicit class weighting guarantees deterministic word alignment and high comma recall.

\begin{table}[h]
\centering
\small
\begin{tabular}{lccc}
\toprule
\textbf{Architecture} & \textbf{Comma F1} & \textbf{Period F1} & \textbf{CPU Latency} \\
\midrule
T5-Small (Seq2Seq) & 0.061 & 0.702 & ~80--140\,ms \\
Afro-XLMR (Ours) & \textbf{0.799} & \textbf{0.934} & \textbf{15--25\,ms} \\
\bottomrule
\end{tabular}
\caption{Comparison of sequence tagging vs. seq2seq.}
\label{tab:ablation}
\end{table}

\section{Edge Deployment \& Quantization Analysis}
To enable accessible deployment across Africa where high-end GPU infrastructure is scarce, we perform post-training dynamic INT8 quantization on the exported ONNX graph. Table~\ref{tab:quantization} compares the precision, latency, and memory footprint between the full half-precision GPU model (\texttt{bfloat16}) and the quantized CPU engine (\texttt{INT8}).

\begin{table}[h]
\centering
\small
\begin{tabular}{lcccc}
\toprule
\textbf{Precision} & \textbf{Comma F1} & \textbf{Period F1} & \textbf{Footprint} & \textbf{Latency} \\
\midrule
\texttt{bfloat16} (GPU) & \textbf{0.799} & \textbf{0.934} & 540\,MB & 3.2\,ms (GPU) \\
\texttt{INT8} (CPU ONNX) & 0.731 & 0.915 & \textbf{278\,MB} & \textbf{15--25\,ms} (CPU) \\
\bottomrule
\end{tabular}
\caption{Comparison of full-precision (\texttt{bfloat16}) vs. dynamic quantized (\texttt{INT8}) ONNX engine.}
\label{tab:quantization}
\end{table}

As shown, INT8 dynamic quantization retains **over 91\% of comma accuracy** and **98\% of sentence-boundary accuracy** while reducing the parameter footprint by 50\% (down to 278\,MB) and delivering sub-30\,ms real-time throughput on standard commodity x86 CPUs.

\paragraph{Zero-Cost Serverless Hosting.} We deploy the INT8 engine using a serverless endpoint on Modal configured to scale to zero after 60 seconds of inactivity, connected to an open static HTML/JS demonstration interface hosted on Hugging Face Spaces.

\section{Conclusion}
In this work, we presented \textbf{AfriPunctCap}, addressing the critical prosodic disconnect between written conventions and spoken pauses across African languages. By performing acoustic pause alignment across 314 languages and fine-tuning a multi-task transformer on multi-sentence streams, we demonstrate state-of-the-art punctuation and capitalization restoration that runs in real time on commodity CPUs. All datasets, models, and code are open-sourced under the AfriSpeech initiative.

\bibliographystyle{IEEEtran}
\bibliography{paper}

\end{document}
